% Options for packages loaded elsewhere
\PassOptionsToPackage{unicode}{hyperref}
\PassOptionsToPackage{hyphens}{url}
\PassOptionsToPackage{space}{xeCJK}
%
\documentclass[
  ignorenonframetext,
  aspectratio=169,
]{beamer}
\title{インプラント FEM：ペリインプラント・バイオフィルムの組成分解力学}
\subtitle{組成 \(\varphi_{Pg}\) → 成長固有ひずみ →
インプラント界面残留応力 → 剥離}
\author{西岡佳祐 --- NIFE / SFB TRR-298}
\date{2026-06-13}

\usepackage{pgfpages}
\setbeamertemplate{caption}[numbered]
\setbeamertemplate{caption label separator}{: }
\setbeamercolor{caption name}{fg=normal text.fg}
\beamertemplatenavigationsymbolsempty
% Prevent slide breaks in the middle of a paragraph
\widowpenalties 1 10000
\raggedbottom
\setbeamertemplate{part page}{
  \centering
  \begin{beamercolorbox}[sep=16pt,center]{part title}
    \usebeamerfont{part title}\insertpart\par
  \end{beamercolorbox}
}
\setbeamertemplate{section page}{
  \centering
  \begin{beamercolorbox}[sep=12pt,center]{part title}
    \usebeamerfont{section title}\insertsection\par
  \end{beamercolorbox}
}
\setbeamertemplate{subsection page}{
  \centering
  \begin{beamercolorbox}[sep=8pt,center]{part title}
    \usebeamerfont{subsection title}\insertsubsection\par
  \end{beamercolorbox}
}
\AtBeginPart{
  \frame{\partpage}
}
\AtBeginSection{
  \ifbibliography
  \else
    \frame{\sectionpage}
  \fi
}
\AtBeginSubsection{
  \frame{\subsectionpage}
}
\usepackage{amsmath,amssymb}
\usepackage{lmodern}
\usepackage{iftex}
\ifPDFTeX
  \usepackage[T1]{fontenc}
  \usepackage[utf8]{inputenc}
  \usepackage{textcomp} % provide euro and other symbols
\else % if luatex or xetex
  \usepackage{unicode-math}
  \defaultfontfeatures{Scale=MatchLowercase}
  \defaultfontfeatures[\rmfamily]{Ligatures=TeX,Scale=1}
  \setmainfont[]{Noto Serif CJK JP}
  \setmonofont[]{Noto Sans Mono CJK JP}
  \ifXeTeX
    \usepackage{xeCJK}
    \setCJKmainfont[]{Noto Serif CJK JP}
  \fi
  \ifLuaTeX
    \usepackage[]{luatexja-fontspec}
    \setmainjfont[]{Noto Serif CJK JP}
  \fi
\fi
\usetheme[]{Madrid}
\usecolortheme{whale}
\usefonttheme{serif} % use mainfont rather than sansfont for slide text
% Use upquote if available, for straight quotes in verbatim environments
\IfFileExists{upquote.sty}{\usepackage{upquote}}{}
\IfFileExists{microtype.sty}{% use microtype if available
  \usepackage[]{microtype}
  \UseMicrotypeSet[protrusion]{basicmath} % disable protrusion for tt fonts
}{}
\makeatletter
\@ifundefined{KOMAClassName}{% if non-KOMA class
  \IfFileExists{parskip.sty}{%
    \usepackage{parskip}
  }{% else
    \setlength{\parindent}{0pt}
    \setlength{\parskip}{6pt plus 2pt minus 1pt}}
}{% if KOMA class
  \KOMAoptions{parskip=half}}
\makeatother
\usepackage{xcolor}
\IfFileExists{xurl.sty}{\usepackage{xurl}}{} % add URL line breaks if available
\IfFileExists{bookmark.sty}{\usepackage{bookmark}}{\usepackage{hyperref}}
\hypersetup{
  pdftitle={インプラント FEM：ペリインプラント・バイオフィルムの組成分解力学},
  pdfauthor={西岡佳祐 --- NIFE / SFB TRR-298},
  hidelinks,
  pdfcreator={LaTeX via pandoc}}
\urlstyle{same} % disable monospaced font for URLs
\newif\ifbibliography
\usepackage{longtable,booktabs,array}
\usepackage{calc} % for calculating minipage widths
\usepackage{caption}
% Make caption package work with longtable
\makeatletter
\def\fnum@table{\tablename~\thetable}
\makeatother
\usepackage{graphicx}
\makeatletter
\def\maxwidth{\ifdim\Gin@nat@width>\linewidth\linewidth\else\Gin@nat@width\fi}
\def\maxheight{\ifdim\Gin@nat@height>\textheight\textheight\else\Gin@nat@height\fi}
\makeatother
% Scale images if necessary, so that they will not overflow the page
% margins by default, and it is still possible to overwrite the defaults
% using explicit options in \includegraphics[width, height, ...]{}
\setkeys{Gin}{width=\maxwidth,height=\maxheight,keepaspectratio}
% Set default figure placement to htbp
\makeatletter
\def\fps@figure{htbp}
\makeatother
\setlength{\emergencystretch}{3em} % prevent overfull lines
\providecommand{\tightlist}{%
  \setlength{\itemsep}{0pt}\setlength{\parskip}{0pt}}
\setcounter{secnumdepth}{-\maxdimen} % remove section numbering
\usepackage{amsmath,amssymb}
\newcommand{\sgn}{\operatorname{sgn}}
\newcommand{\relu}[1]{\left[#1\right]_{+}}
\ifLuaTeX
  \usepackage{selnolig}  % disable illegal ligatures
\fi

\begin{document}
\frame{\titlepage}

\begin{frame}{本デッキの位置づけ}
\protect\hypertarget{ux672cux30c7ux30c3ux30adux306eux4f4dux7f6eux3065ux3051}{}
\textbf{Implant FEM} ---
ペリインプラント・バイオフィルムの力学を組成分解で扱う。

データの流れ：

raw 16S → guild \(\varphi\) → gLV/Hamilton（＋符号 prior）→ 空間 PDE →
\textbf{FEM 力学}

体制（supervisor context）：

\begin{itemize}
\tightlist
\item
  \textbf{Soleimani}（IKM, FEM）--- 連続体力学・有限要素の枠組み
\item
  \textbf{Junker}（IKM, Hamilton 原理）--- 材料モデルの変分定式化
\item
  \textbf{Szafrański}（MHH / NIFE）--- 口腔バイオフィルム・臨床
\end{itemize}

\vspace{0.4em}

\textbf{一行命題：} dysbiosis を
\emph{力学的な}インプラント界面リスクとして読み出す。
\end{frame}

\begin{frame}{臨床的問題：ペリインプラント炎}
\protect\hypertarget{ux81e8ux5e8aux7684ux554fux984cux30daux30eaux30a4ux30f3ux30d7ux30e9ux30f3ux30c8ux708e}{}
ペリインプラント炎はバイオフィルム駆動の炎症であり、インプラント周囲の
\textbf{支持骨喪失}を引き起こす。

\begin{itemize}
\tightlist
\item
  インプラント患者の \textbf{20--35\%} に発症 ---
  インプラント失敗の主因。
\item
  \textbf{NIFE} =
  インプラント研究開発拠点。基材（substratum）は\textbf{チタン}。
\item
  \textbf{Dieckow コホート} =
  アバットメント上のバイオフィルムを縦断観察。
\end{itemize}

\vspace{0.4em}

力学的に何が起きているか：バイオフィルムが Ti 界面で成長し、界面に
\textbf{応力}を生み、やがて\textbf{剥離}する。これを連続体力学で定量化する。
\end{frame}

\begin{frame}{アイデア：組成から力学への連鎖}
\protect\hypertarget{ux30a2ux30a4ux30c7ux30a2ux7d44ux6210ux304bux3089ux529bux5b66ux3078ux306eux9023ux9396}{}
中心となる因果連鎖：

\[\varphi_{Pg}\;\rightarrow\;\text{成長固有ひずみ}\;\rightarrow\;\text{界面残留応力}\;\rightarrow\;\text{剥離}\;\rightarrow\;\text{ペリインプラント炎}\]

\includegraphics[width=\textwidth,height=0.46\textheight]{masterarbeit_ansys_fem/figures/fem_clinical_schematic.pdf}

CH（常在）対 DH（dysbiotic）バイオフィルム下の Ti
ねじ山：成長誘起の界面応力と剥離。
\end{frame}

\begin{frame}{方法：Hamilton 原理の材料モデル}
\protect\hypertarget{ux65b9ux6cd5hamilton-ux539fux7406ux306eux6750ux6599ux30e2ux30c7ux30eb}{}
各 Gauss 点で\textbf{現象論則の代わりに} Python の Hamilton
原理材料モデルを呼ぶ。

\begin{itemize}
\tightlist
\item
  乗法的な成長分解 \(F = F_e\,F_g\) ＋ neo-Hookean 弾性。
\item
  成長 \(F_g\) は組成 \(\varphi_{Pg}\) の固有ひずみで駆動。
\item
  弾性応答 \(F_e\) が応力を担い、界面で残留応力として蓄積。
\end{itemize}

\vspace{0.3em}

\textbf{検証（Abaqus）：}

\begin{itemize}
\tightlist
\item
  微小ひずみ：\(\sigma = E\,\varepsilon\) と一致。
\item
  有限ひずみ（NLGEOM）：neo-Hookean 解と一致。
\end{itemize}

\vspace{0.3em}

\(\Rightarrow\) 材料則は外部 FEM
ソルバへそのまま差し込める（カーネルとして検証済み）。
\end{frame}

\begin{frame}{成長キャリブレーション}
\protect\hypertarget{ux6210ux9577ux30adux30e3ux30eaux30d6ux30ecux30fcux30b7ux30e7ux30f3}{}
成長固有ひずみを \textbf{CLSM 深さプロファイル}からキャリブレーション。

\includegraphics[width=\textwidth,height=0.52\textheight]{masterarbeit_ansys_fem/figures/F2_growth_calibration.pdf}

\begin{itemize}
\tightlist
\item
  DH：深さ相関 \(0.88\)、\(\beta_{\text{depth}} \approx 17.2\)。
\item
  常在（commensal）コントロールは\textbf{逆符号} ---
  単なる体積膨張ではない。
\end{itemize}
\end{frame}

\begin{frame}{結果：残留応力カラム \(S_{11}(z)\)}
\protect\hypertarget{ux7d50ux679cux6b8bux7559ux5fdcux529bux30abux30e9ux30e0-s_11z}{}
\includegraphics[width=\textwidth,height=0.54\textheight]{masterarbeit_ansys_fem/figures/F4_residual_stress_DHvsCH.pdf}

\begin{itemize}
\tightlist
\item
  応力は\textbf{基材}（インプラント--アバットメント壁）に集中、自由な上端
  \(\approx 0\)。
\item
  \(S_{11}(z)\) は \(\varphi_{Pg}(z)\) を\textbf{追従} --- 深部 Pg
  負荷が界面応力を作る。
\end{itemize}
\end{frame}

\begin{frame}{結果：剥離（cohesive 界面）}
\protect\hypertarget{ux7d50ux679cux5265ux96e2cohesive-ux754cux9762}{}
\includegraphics[width=\textwidth,height=0.52\textheight]{masterarbeit_ansys_fem/figures/F6_delamination_DHvsCH.pdf}

同一の界面強度で cohesive-interface FEM を回す：

\begin{itemize}
\tightlist
\item
  \textbf{dysbiotic 83\%} 対 \textbf{commensal 25\%}
  が剥離（\textbf{3.3×}）。
\item
  唯一の原因 = \textbf{Pg 深さプロファイル}（強度・荷重は同一）。
\end{itemize}
\end{frame}

\begin{frame}{3D ヒーロー：ペリインプラント全アセンブリ}
\protect\hypertarget{d-ux30d2ux30fcux30edux30fcux30daux30eaux30a4ux30f3ux30d7ux30e9ux30f3ux30c8ux5168ux30a2ux30bbux30f3ux30d6ux30ea}{}
\includegraphics[width=\textwidth,height=0.5\textheight]{masterarbeit_ansys_fem/figures/fem_implant_screw3d.pdf}

インプラント＋隣接歯＋共有歯槽骨。咬合荷重は\textbf{歯頸部のペリインプラント骨}に
集中（marginal-bone-loss
シグネチャ）。組成メカニズムは\textbf{現実の解剖学へ そのまま転移}する。
\end{frame}

\begin{frame}{クラウン荷重伝達：モーメントアーム効果}
\protect\hypertarget{ux30afux30e9ux30a6ux30f3ux8377ux91cdux4f1dux9054ux30e2ux30fcux30e1ux30f3ux30c8ux30a2ux30fcux30e0ux52b9ux679c}{}
\includegraphics[width=\textwidth,height=0.5\textheight]{masterarbeit_ansys_fem/figures/fem_implant_crown_fem.pdf}

実歯クラウンをアバットメントに装着し、荷重を\textbf{咬合面（z≈38,
隣在歯と同一平面）}
で受けると、首部まわりの\textbf{曲げモーメントが約3.4倍}に。生理的咬合（100
N, 30°, ISO 14801）でも\textbf{ペリインプラント歯槽頂骨のピーク応力が
×1.6（55→88 MPa）}。 \(\Rightarrow\)
補綴設計（クラウン高・荷重位置）が\textbf{辺縁骨リスクに直接効く}。
\end{frame}

\begin{frame}{2D 連成セクション（補助）}
\protect\hypertarget{d-ux9023ux6210ux30bbux30afux30b7ux30e7ux30f3ux88dcux52a9}{}
\includegraphics[width=\textwidth,height=0.52\textheight]{masterarbeit_ansys_fem/figures/fem_tier2b_real.pdf}

連成 implant＋tooth＋bone の近遠心断面。(A) 多材料解剖、(B) 咬合 von
Mises の 共有骨連成 --- \textbf{歯頸部応力 =
ペリインプラント辺縁骨喪失のシグネチャ}。
\end{frame}

\begin{frame}{患者ブリッジ}
\protect\hypertarget{ux60a3ux8005ux30d6ux30eaux30c3ux30b8}{}
\includegraphics[width=\textwidth,height=0.48\textheight]{masterarbeit_ansys_fem/figures/F1c_patient_bridge_validation.pdf}

Dieckow 10 患者を検証済みの剥離カーネル \(J\) へ射影：

\begin{itemize}
\tightlist
\item
  患者間ばらつき \textbf{\(CV = 0.86\)} は\textbf{測定された} Pg
  負荷が決める。
\item
  借用した in-vitro 深さ形状は \textbf{34× 弱いレバー}（\(\pm 8\%\)）。
\end{itemize}

\vspace{0.2em}

\(\Rightarrow\) \textbf{0-D の臨床 Pg
読み値}が順序的な剥離リスクの代理として十分。
\end{frame}

\begin{frame}{検証・限界・展望}
\protect\hypertarget{ux691cux8a3cux9650ux754cux5c55ux671b}{}
\textbf{正直な評価：}

\begin{itemize}
\tightlist
\item
  成長の\textbf{形状}は CLSM から較正済み。
\item
  応力\textbf{絶対値}はパラメトリック（バイオフィルム弾性率は約 7 桁の幅
  → \(\sigma/\mu\) で報告）。
\item
  主張は\textbf{順序的・機構的}であり予測的ではない。差別化は\emph{バイオフィルム連成}
  であって汎用インプラント力学ではない。
\end{itemize}

\textbf{検証実験：} 深さ分解 AFM／剪断流剥離 対 Pg 負荷。

\vspace{0.3em}

\textbf{展望：}
成長固有ひずみ→界面応力→剥離の能力は、慶應・村松研の計算固体力学
および半導体薄膜応力・ダイシングストリート剥離（DISCO）へ転移する。
\end{frame}

\begin{frame}[fragile]{付録：FEM 条件（材料・境界・荷重）}
\protect\hypertarget{ux4ed8ux9332fem-ux6761ux4ef6ux6750ux6599ux5883ux754cux8377ux91cd}{}
\begin{longtable}[]{@{}lrr@{}}
\toprule
材料 & \(E\) (MPa) & \(\nu\) \\
\midrule
\endhead
チタン (Ti-6Al-4V) & 110 000 & 0.34 \\
皮質骨 / 海綿骨 & 13 700 / 1 000 & 0.30 \\
象牙質 & 18 000 & 0.31 \\
PDL（0.25 mm, 線形理想化） & 50 & 0.45 \\
クラウン（リチウムジシリケート / ジルコニア） & 95 000 / 210 000 &
0.30 \\
歯肉（粘膜） & 3 & 0.45 \\
バイオフィルム（＋成長 \(\varepsilon_g=0.19\)） & 1.0 & 0.45 \\
\bottomrule
\end{longtable}

\textbf{境界}：骨の crop
面で完全固定（\(U_{1,2,3}=0\)）。\textbf{界面}：オッセオ＝*TIE、クラウン＝アバットメント天面に
*TIE、脱統合変種のみ接触 \(\mu=0.3\)。 \textbf{荷重}：ISO 14801、100
N・30°斜め（lat/ax=0.577）。Step1=成長固有ひずみ、Step2=咬合（クラウンありは\textbf{咬合面
z=38}＝モーメントアーム）。 \textbf{要素}：連成=C3D4（〜27万）、ISO14801
クーポン=C3D10（線形は集中を 9--15\% 過小）、収束 ±7\%。単位
mm-N-MPa。詳細は \texttt{FEM\_CONDITIONS.md}。
\end{frame}

\end{document}
